\paragraph{Leakage.}
With a random, sample-level split of the same size, in which every signer appears in both training
and test data, the Conv-Transformer reaches \numLeakRandom{}\,\% test top-1 instead of
\numLeakSigner{}\,\%. A random split thus overestimates the accuracy on new users by \numLeakGap{}
points, more than the effect of any architectural or preprocessing choice studied here.

\paragraph{Ablations.}
\Cref{fig:ablations} changes one factor at a time with respect to the full model
(\numRefValTop{}\,\% validation top-1; seed-to-seed standard deviation of 0.1 point). Early stopping
proved to be a confounding factor. Two variants initially stopped on a validation plateau after 21
and 31 epochs, while the learning rate was still high, and appeared to lose 4.4 and 2.7 points.
Re-run over the full 50-epoch schedule, their effects reduce to \numAblNoHandlocal{} (no hand-local
features) and \numAblXyz{} (depth added) points. The other variants marked with $\dagger$
(\numAblEarlyStopped{} in total) stopped between epochs 41 and 46; their effects may be slightly
overestimated, and our compute budget did not allow re-running them.

